(1) 能量高于台阶时 $k_1=\sqrt{2mE}/\hbar$、$k_2=\sqrt{2m(E-V_0)}/\hbar$。取 $\psi_{x<0}=e^{ik_1x}+r e^{-ik_1x}$、$\psi_{x>0}=t e^{ik_2x}$；在 $x=0$ 处 $1+r=t$、$k_1(1-r)=k_2t$，得 $r=(k_1-k_2)/(k_1+k_2)$、$t=2k_1/(k_1+k_2)$。
(2) 入射与反射粒子流分别为 $\hbar k_1/m$ 与 $-\hbar k_1|r|^2/m$，右侧透射流为 $\hbar k_2|t|^2/m$。因此 $R=|r|^2=[(k_1-k_2)/(k_1+k_2)]^2$，$T=(k_2/k_1)|t|^2=4k_1k_2/(k_1+k_2)^2$；展开分子立即得 $R+T=1$，不能只把振幅平方当透射率。
(3) 若 $E<V_0$，令 $\kappa=\sqrt{2m(V_0-E)}/\hbar$，无穷远有界解是 $\psi_{x>0}=t e^{-\kappa x}$。它可在阶跃右侧留下非零概率密度，但此纯实指数的概率流 $j=(\hbar/m)\Im(\psi^*\psi')=0$；稳态全反射 $R=1$，无无限远右行传播通量。
(4) $V_0\to0$ 时 $k_2\to k_1$、$R\to0,T\to1$；由上趋近 $E\downarrow V_0$ 时 $k_2\to0^+$，得 $T\to0,R\to1$。无限延伸的势阶跃不同于有限宽度势垒，后者右端恢复传播区，倏逝解可以重新连接成非零透射流，即量子隧穿。
